%%%% LISTA DE ABREVIATURAS, SIGLAS E ACRÔNIMOS

\listadeabrevsiglaseacr%% Formatação da lista de abreviaturas, siglas e acrônimos

\begin{Spacing}{1.5}
\section*{\nomesiglas}
\begin{tabularx}{\textwidth}{@{}lX@{}}
ABNT  & Associação Brasileira de Normas Técnicas \\
API   & Interface de Programação de Aplicações, do inglês \textit{Application Programming Interface} \\
ARM   & Família de arquiteturas de processadores RISC utilizada no subsistema embarcado \\
CPU   & Unidade Central de Processamento, do inglês \textit{Central Processing Unit} \\
DMA   & Acesso Direto à Memória, do inglês \textit{Direct Memory Access} \\
FAST  & FIX Adapted for STreaming \\
FIX   & Financial Information eXchange \\
FPGA  & Arranjo de Portas Programáveis em Campo, do inglês \textit{Field-Programmable Gate Array} \\
GHDL  & Simulador de VHDL utilizado nos testes automatizados \\
HDL   & Linguagem de Descrição de Hardware, do inglês \textit{Hardware Description Language} \\
HFT   & Negociação de Alta Frequência, do inglês \textit{High-Frequency Trading} \\
HLS   & Síntese de Alto Nível, do inglês \textit{High-Level Synthesis} \\
HPS   & Sistema Processador Rígido, do inglês \textit{Hard Processor System} \\
IP    & Propriedade Intelectual reutilizável de hardware, do inglês \textit{Intellectual Property} \\
MMIO  & Entrada e Saída Mapeada em Memória, do inglês \textit{Memory-Mapped Input/Output} \\
PDF   & Formato de Documento Portátil, do inglês \textit{Portable Document Format} \\
RTL   & Nível de Transferência entre Registradores, do inglês \textit{Register Transfer Level} \\
SoC   & Sistema em um Chip, do inglês \textit{System-on-Chip} \\
TCP   & Protocolo de Controle de Transmissão, do inglês \textit{Transmission Control Protocol} \\
UTFPR & Universidade Tecnológica Federal do Paraná \\
VCD   & Value Change Dump \\
VHDL  & VHSIC Hardware Description Language \\
\end{tabularx}
\end{Spacing}
